% ─────────────────────────────────────────────────────────────────────────────
\section{Related Work}
\label{sec:related}

\paragraph{Cross-view retrieval for UAVs.}
Most UAV geo-localization systems cast the problem as retrieval of a
geo-tagged satellite tile. Transformer and part-based embeddings
(TransGeo~\cite{zhu2022transgeo}, CAMP~\cite{wu2024camp}) are strong on
aligned drone--satellite pairs. Newer benchmarks remove the alignment
assumption: GTA-UAV/Game4Loc~\cite{ji2025game4loc} introduces partial overlap
and metric localization over a continuous area, and
World-UAV/UAV-GeoLoc~\cite{wu2025uavgeoloc} a one-to-$N$ setting with large
rotation and scale discrepancies, under which global descriptors degrade
markedly. Retrieval returns a tile, not a position; a metric estimate
requires a second stage. Our approach is complementary: it does not learn an
appearance embedding, and it is equivariant to rotation and scale by
construction because both are explicit search variables.

\paragraph{Matching, registration and pose.}
The second family matches local or dense features between the UAV frame and
a reference crop and solves for a homography or a PnP pose.
LightGlue~\cite{lindenberger2023lightglue} and RoMa~\cite{edstedt2024roma}
are representative matchers; AirLock+~\cite{deng2026airlock} registers
UAV video to satellite imagery with homographies over long real trajectories,
and AnyVisLoc~\cite{ye2026anyvisloc} standardizes a
retrieval$\rightarrow$matching$\rightarrow$PnP pipeline for low-altitude,
multi-view flight with aerial and satellite 2.5D references.
Bearing-UAV~\cite{liu2026bearinguav} jointly estimates position and heading
from unaligned views. These pipelines rely on texture correspondences; our
evidence is instead the building layout, and our registration is a global,
exhaustive optimization over a pose group rather than a consensus over
putative correspondences~\cite{fischler1981ransac}.

\paragraph{Explicit structure and vector maps.}
Building footprints have been used as a map primitive by descriptor methods
such as BRM~\cite{choi2020brm}, by coarse-to-fine vector-map positioning with
a particle filter (CFBVM)~\cite{ouyang2024cfbvm}, and by shape-and-relationship
matching on clean vector e-maps~\cite{ssrm2024}. These methods either require
an external vector layer or integrate evidence over a trajectory. The
conference version of this work encoded individual CDT triangles with
Ekeland free-cone angles and retrieved them from a raster-derived index; its
failure analysis (Sec.~\ref{sec:why_m0}) motivates the present
configuration-level formulation.

\paragraph{Chamfer and correlation-based registration.}
Our objective belongs to a long line of template-based registration:
chamfer matching~\cite{barrow1977chamfer,borgefors1988hcma}, FFT-based
registration over translation, rotation and scale~\cite{reddy1996fft}, and
fast normalized cross-correlation~\cite{lewis1995ncc}, including its masked
form for partially overlapping or irregular supports. We do not claim these
tools as novel. The contribution is how they are assembled and characterized
for GNSS-denied UAV localization from \emph{raster-derived} structure: the
uncertainty of the structure extraction is propagated into the evidence,
the pose uncertainty is quantified, a building-shape verifier resolves
aliasing, and an integrity model decides when the structural evidence is
insufficient---together with an explicit measurement of the regime in which
footprint geometry is observable at all.

\paragraph{Uncertainty and integrity.}
Uncertainty in cross-view localization has mostly been modeled for training
correspondences or as learned confidence heads. We instead derive confidence
at inference time from observable geometric evidence---peak distinctiveness,
the curvature of the objective (a Laplace approximation~\cite{mackay2003information}),
the number, spread and persistence of observed buildings---and calibrate it
on held-out sites.
